\section{Principle VII: Property}

Property is the natural safeguard of man, an industrious animal still more than reasonable, for whom his needs, his weakness, the dissatisfaction that he brings from the cradle make a narrow duty of transforming what surrounds him. Pascal is wrong to scoff at those poor children who say: this coin is mine.

One would not say ``me", without saying ``mine". Without property, man is condemned to death.

To possess is to command, it is to dispose of oneself, it is to be able to resist others, it is to exercise an influence, were it only by reaction.

Property emancipates existence and confers an authority at least on the goods of the earth and the fruits of labor.



\flrightit{Posted on 2012-12-02 by Cologero }

\begin{center}* * *\end{center}

\begin{footnotesize}\begin{sffamily}



\texttt{Jason-Adam on 2012-12-03 at 17:57 said: }

Very direct and to the point. 

It does create a dilemma for those ``Traditionalists" so-called who see hope in communism.


\hfill

\texttt{Saladin on 2012-12-03 at 19:14 said: }

lol! ``Traditionalists who see hope in communism"!!! You must be joking, right?


\hfill

\texttt{Jason-Adam on 2012-12-03 at 19:34 said: }

I wish I were but I am not….

There are men who claim to be both followers of Rene Guenon and Karl Marx. Alexandr Dugin's writings are a prime example of this tendency…..

The person who first introduced me to Baron Evola's books was a National Communist…..


\hfill

\texttt{Jason-Adam on 2012-12-03 at 19:37 said: }

Read this to see how Dugin tries to mix Evolian-ism with Marxism.

\url{http://against-postmodern.org/section-iv-mission-julius-evola}

I can only imagine what the Baron would say if he could read this…..

More and more, I am starting to see Dugin as the prophecied counter-intiation that Guenon mentioned in Reign of Quantity….


\hfill

\texttt{Cologero on 2012-12-03 at 21:02 said: }

It seems that a lingering affection for communism is the equivalent to the Stockholm syndrome\footnote{\url{http://en.wikipedia.org/wiki/Stockholm_syndrome}} for some Russian intellectuals. Another we can mention is Israel Shamir.

Without property, man is a slave to the state. Even Medieval peasants had their land and couldn't be evicted by bankers or tax collectors. A Romanian friend has told me tales about life under communism. For example, the authorities deemed that his grandmother had more house than she needed. So they took over some of its rooms and moved in another family and a retired military man. She had no say in the matter.


\hfill

\texttt{Dominion on 2012-12-04 at 00:43 said: }

The right to property is fundamental in order to preserve the ability to attain freedom (the Traditionalist notes, of course, that it is not freedom in and of itself, as the American right would have us believe). However, the right is not absolute. The banker or corporate shark who exploits the worker and nation cannot hide behind his right to property when the State exercises its own rightful duty to uphold Justice.

On the topic of communism, two important points must be made: first, that ``communism" is many times used as an umbrella term to describe varying and contradictory elements, making it very much the equivalent of the ``fascist" epithet. Second, that the way in which certain Traditionalists view ``socialism" (that is to say, a society in which the worker has some sort of democratic control over capital and where the State advances these reforms) as being far from antithetical to the aims of Tradition is far more nuanced than some seem to think. The best criticism I have heard from these branches is that the exploited worker is viewed as a mere unit of production; socialism's aim is to view the worker as a human being whose work cannot be separated from his inner development. Thus, a socialism which allows workers democratic control over capital is necessary in order to allow the worker to develop his various mental and spiritual capacities. 

With regards to the exoteric institutions of Tradition, many movements today which can be called ``socialist" have moved in this direction, or at least have the potential to be educated towards this framework. One example is the Russian Communist Party today, which supports the Orthodox Church being a foundation stone of Russia, and aims to maintain social systems which view the family as a social good and unit, such as the public kindergarten system. It also wants to maintain small and medium sized businesses, which is not so distant from the distributist view of property's role. Dugin's fourth political theory is useful in that, using Being as its basis, it takes what is useful from the previous theories, including socialism. But what is base must be made subtle; what is made from the dust of the so-called enlightenment must be baptized in the water which brings true enlightenment. Be it property and the market, socialism and the worker, or the fasces and the authority of the State, this will need to be done in order to build an order in which the values of Tradition are embodied.


\hfill

\texttt{Jason-Adam on 2012-12-04 at 02:54 said: }

Cologero is correct on the need for all persons to be property owners. Property is the basis of an individual life as Maurras articulately said above…..

I often wondered myself why Russians like Dugin who were victims of Communist repression suddenly turned nostalgic for the USSR….

Do you think though there's any merit in the arguement that since the USSR fell before the USA it was therefore not as degenerated ?


\hfill

\texttt{Saladin on 2012-12-04 at 19:59 said: }

I am only vaguely familiar with Dugin. His Geopolitics is certainly interesting. Tradition has no use for Communism, and communism is totally opposed to Tradition. As Cologero said ``without property, man is a slave to the state." There has never been a Tradition that has not protected the concept of Private Property.


\hfill

\texttt{Saladin on 2012-12-04 at 20:06 said: }

Thanks for the link Jason! Dugin suggests that we read Evola ``from the Left"! That explains a lot! He clearly has his own reading of Tradition but I completely agree with Cologero that this may be a bit of Stockholm Syndrome and definitely not ``the prophecied counter-intiation" figure.


\end{sffamily}\end{footnotesize}
